\SourcePage{113}{118}
\section{Charge conjugation and time reversal of spinors}

The factors $\psi_{p\sigma} = u_{p\sigma}e^{-ipx}$ that multiply the operators $\widehat a_{\mathbf{p}\sigma}$ in~(25.1) are the wave functions of free particles (which we shall call ``electrons'') with momenta $\mathbf{p}$ and polarizations $\sigma$:
\[
\psi^{(\text{e})}_{p\sigma} = \psi_{p\sigma}.
\]

The factors $\bar\psi_{-p\,-\sigma}$ multiplying the operators $\widehat b_{\mathbf p\sigma}$ must instead be regarded as the wave functions of positrons with the same $\mathbf{p}$ and $\sigma$. It will turn out, however, that the electron and positron functions are expressed in different bispinor representations. This follows from the different transformation properties of $\psi$ and $\bar\psi$ and from the different systems of equations satisfied by their components. To remove this defect, the components of $\bar\psi_{-p\,-\sigma}$ must undergo a certain unitary \SourcePage{114}{119}transformation such that the new four-component function satisfies the same equation as $\psi_{p\sigma}$\,\footnote{For particles of spin~0 this issue did not arise at all: the scalar functions $\psi$ and $\psi^*$ satisfy the same equation, and $\psi^*_{-p}$ simply coincides with $\psi_p$.}. This is the function that we shall call the positron wave function, with momentum $\mathbf{p}$ and polarization $\sigma$. Denoting the required unitary-transformation matrix by $U_C$, we write
\begin{equation}
\psi^{(\text{p})}_{p\sigma} = U_C\widetilde{\bar\psi}_{-p\,-\sigma}.
\tag{26.1}
\end{equation}

The operation $C$ that produces this function from $\psi_{-p\,-\sigma}$ is called \emph{charge conjugation} of the wave function (\emph{H.\,A.~Kramers}, 1937). The concept is of course not restricted to plane waves. For any function $\psi$ there is a ``charge-conjugate'' function
\begin{equation}
\widehat C\psi(t,\mathbf{r}) = U_C\widetilde{\bar\psi}(t,\mathbf{r}),
\tag{26.2}
\end{equation}
which transforms like $\psi$ and satisfies the same equation.

The properties of $U_C$ follow from this definition. If $\psi$ is a solution of the Dirac equation $(\gamma\widehat p - m)\psi = 0$, then $\bar\psi$ satisfies
\[
\bar\psi(\gamma\widehat p + m) = 0,\qquad \text{or}\quad (\widetilde\gamma\widehat p + m)\widetilde{\bar\psi} = 0.
\]
Multiplication on the left by $U_C$ gives
\[
U_C\widetilde\gamma\widehat p\,\widetilde{\bar\psi} + mU_C\widetilde{\bar\psi} = 0,
\]
We require $U_C\widetilde{\bar\psi}$ to satisfy the same equation as $\psi$:
\[
(\gamma\widehat p - m)U_C\widetilde{\bar\psi} = 0.
\]
Comparison gives the following ``commutation relation'' between $U_C$ and the matrices $\gamma^\mu$\,\footnote{We also note the resulting equality
\begin{equation}
U_C\widetilde\gamma^5 = \gamma^5U_C.
\tag{26.3a}
\end{equation}}:
\begin{equation}
U_C\widetilde\gamma^\mu = -\gamma^\mu U_C.
\tag{26.3}
\end{equation}

In what follows, we assume that the wave functions are given in the spinor or standard representation; the general case of an arbitrary representation will be considered only at the end of this section. In these representations,
\begin{equation}
\begin{aligned}
\gamma^{0,2} &= \widetilde\gamma^{0,2},&\qquad
\gamma^{1,3} &= -\widetilde\gamma^{1,3},\\
\bigl(\gamma^{0,1,3}\bigr)^* &= \gamma^{0,1,3},&
\gamma^{2*} &= -\gamma^2.
\end{aligned}
\tag{26.4}
\end{equation}
The conditions~(26.3) are then satisfied by $U_C = \eta_C\gamma^2\gamma^0$, where $\eta_C$ is an arbitrary constant. The requirement $\widehat C^2 = 1$ gives $|\eta_C|^2 = 1$, so $U_C$ is fixed only up to a \SourcePage{115}{120}phase factor. We shall choose $\eta_C = 1$ from now on, giving
\begin{equation}
U_C = \gamma^2\gamma^0 = -\alpha_y.
\tag{26.5}
\end{equation}
Since also $\bar\psi = \psi^*\gamma^0 = \widetilde\gamma^0\psi^* = \gamma^0\psi^*$, the action of $\widehat C$ can be written
\begin{equation}
\widehat C\psi = \gamma^2\gamma^0\widetilde{\bar\psi} = \gamma^2\psi^*.
\tag{26.6}
\end{equation}

Explicitly, transformation~(26.6) in the spinor representation is
\begin{equation}
C:\quad \xi^\alpha \to -i\eta^{\dot\alpha*},\qquad
\eta_{\dot\alpha} \to -i\xi^*_\alpha.
\tag{26.7a}
\end{equation}
Equivalently,
\begin{equation}
C:\quad \xi_\alpha \to -i\eta^*_{\dot\alpha},\qquad
\eta^{\dot\alpha} \to -i\xi^{\alpha*}.
\tag{26.7b}
\end{equation}

Charge conjugation of the plane waves $\psi_{\pm p\sigma}$ is readily performed using their explicit expressions~(23.9) and the standard-representation matrix
\begin{equation}
U_C = \begin{pmatrix} 0 & -\sigma_y\\ -\sigma_y & 0 \end{pmatrix}.
\tag{26.8}
\end{equation}
Noting that
\[
\sigma_y\sigma^* = -\boldsymbol{\sigma}\sigma_y,
\]
and defining $w^{(\sigma)\prime}$ as in~(23.16), we obtain
\begin{equation}
U_C\widetilde{\bar u}_{-p\,-\sigma} = u_{p\sigma},\qquad
U_C u_{-p\,-\sigma} = \widetilde{\bar u}_{p\sigma}.
\tag{26.9}
\end{equation}
Thus
\begin{equation}
\widehat C\psi_{-p\,-\sigma} = \psi_{p\sigma},
\tag{26.10}
\end{equation}
so the functions $\psi_{-p\,-\sigma}$ that occur in the $\psi$-operators~(25.1) together with $\widehat b_{p\sigma}$ do indeed correspond to particle states of momentum $\mathbf{p}$ and polarization $\sigma$. We also see that electron and positron states are described by the same functions:
\[
\psi^{(\text{e})}_{p\sigma} = \psi^{(\text{p})}_{p\sigma} = \psi_{p\sigma}.
\]

This is natural, since the functions $\psi_{p\sigma}$ contain information only about the particle's momentum and polarization.

Time reversal may be considered in an analogous way. Reversing the sign of time must be accompanied by complex conjugation of the wave function. To obtain the resulting time-reversed wave function $(\widehat T\psi)$ in the same representation as the original $\psi$, the components of $\psi^*$ (or $\bar\psi$) must \SourcePage{116}{121}also undergo a unitary transformation. In analogy with~(26.2), we therefore represent the action of $\widehat T$ on $\psi$ as
\begin{equation}
\widehat T\psi(t,\mathbf{r}) = U_T\widetilde{\bar\psi}(-t,\mathbf{r}),
\tag{26.11}
\end{equation}
where $U_T$ is unitary.

We again write the Dirac equation satisfied by $\psi$:
\[
\Bigl(i\gamma^0\frac{\partial}{\partial t} + i\boldsymbol{\gamma}\boldsymbol{\nabla} - m\Bigr)\psi(t,\mathbf{r}) = 0,
\]
and the equation for $\bar\psi$:
\[
\Bigl(i\widetilde\gamma^0\frac{\partial}{\partial t} + i\widetilde{\boldsymbol{\gamma}}\boldsymbol{\nabla} + m\Bigr)\widetilde{\bar\psi}(t,\mathbf{r}) = 0.
\]
In the latter equation, replace $t$ by $-t$ and multiply from the left by $-U_T$:
\[
\Bigl(iU_T\widetilde\gamma^0\frac{\partial}{\partial t} - iU_T\widetilde{\boldsymbol{\gamma}}\boldsymbol{\nabla}\Bigr)\widetilde{\bar\psi}(-t,\mathbf{r}) - mU_T\widetilde{\bar\psi}(-t,\mathbf{r}) = 0.
\]
We require $U_T\widetilde{\bar\psi}(-t,\mathbf{r})$ to satisfy the same equation as $\psi(t,\mathbf{r})$:
\[
\Bigl(i\gamma^0\frac{\partial}{\partial t} + i\boldsymbol{\gamma}\boldsymbol{\nabla}\Bigr)U_T\widetilde{\bar\psi}(-t,\mathbf{r}) - mU_T\widetilde{\bar\psi}(-t,\mathbf{r}) = 0.
\]
Comparison of the two equations shows that $U_T$ must obey
\begin{equation}
U_T\widetilde\gamma^0 = \gamma^0 U_T,\qquad U_T\widetilde{\boldsymbol{\gamma}} = -\boldsymbol{\gamma}U_T.
\tag{26.12}
\end{equation}
In the spinor and standard representations, these conditions are satisfied by\,\footnote{The choice of phase factor in~(26.13) is related to that in~(26.5) for the reasons stated below in the note on p.~124.}
\begin{equation}
U_T = i\gamma^3\gamma^1\gamma^0.
\tag{26.13}
\end{equation}

The action of $\widehat T$ is consequently
\begin{equation}
\widehat T\psi(t,\mathbf{r}) = i\gamma^3\gamma^1\gamma^0\widetilde{\bar\psi}(-t,\mathbf{r}) = i\gamma^3\gamma^1\psi^*(-t,\mathbf{r}).
\tag{26.14}
\end{equation}
Explicitly, this transformation in the spinor representation is
\begin{equation}
T:\quad \xi^\alpha \to -i\xi^*_\alpha,\qquad \eta_{\dot\alpha} \to i\eta^{\dot\alpha*}
\tag{26.15a}
\end{equation}
or
\begin{equation}
T:\quad \xi_\alpha \to i\xi^{\alpha*},\qquad \eta^{\dot\alpha} \to -i\eta^*_{\dot\alpha}.
\tag{26.15b}
\end{equation}

In the standard representation,
\begin{equation}
U_T = \begin{pmatrix} \sigma_y & 0\\ 0 & -\sigma_y \end{pmatrix}.
\tag{26.16}
\end{equation}

\SourcePage{117}{122}
We now find the result of applying all three operations $P$, $T$, and $C$ to $\psi$. Successively,
\[
\widehat T\psi(t,\mathbf{r}) = -i\gamma^1\gamma^3\psi^*(-t,\mathbf{r}),
\]
\[
\widehat P\widehat T\psi(t,\mathbf{r}) = i\gamma^0(\widehat T\psi) = \gamma^0\gamma^1\gamma^3\psi^*(-t,-\mathbf{r}),
\]
\[
\widehat C\widehat P\widehat T\psi(t,\mathbf{r}) = \gamma^2(\gamma^0\gamma^1\gamma^3\psi^*)^* = \gamma^2\gamma^0\gamma^1\gamma^3\psi(-t,-\mathbf{r}),
\]
or
\begin{equation}
\widehat C\widehat P\widehat T\psi(t,\mathbf{r}) = i\gamma^5\psi(-t,-\mathbf{r}).
\tag{26.17}
\end{equation}

In the spinor representation,
\begin{equation}
CPT:\quad \xi^\alpha \to -i\xi^\alpha,\qquad \eta_{\dot\alpha} \to i\eta_{\dot\alpha},
\tag{26.18}
\end{equation}
as can also be checked directly from the transformation rules~(20.4), (26.7), and (26.15)\,\footnote{The notation $\widehat C\widehat P\widehat T$ means that the operators act from right to left. The overall sign in~(26.17), (26.18) depends on their order because $\widehat T$ does not commute with $\widehat C$ or $\widehat P$ in their action on a bispinor.}.

The expressions above for $U_C$ and $U_T$ assumed the spinor or standard representation of $\psi$. We finally determine which of their properties persist in an arbitrary representation.

If $\psi$ undergoes a unitary transformation,
\begin{equation}
\psi' = U\psi,\quad \gamma' = U\gamma U^{-1},\quad \bar\psi' = \psi'^*\gamma^{0\prime} = \bar\psi U^+ = \bar\psi\widetilde U^{-1},
\tag{26.19}
\end{equation}
then, in the new representation,
\[
(\widehat C\psi)' = U(\widehat C\psi) = UU_C\widetilde{\bar\psi} = UU_C(\widetilde{\bar\psi'U}) = UU_C\widetilde U\widetilde{\bar\psi'}.
\]
Comparing this with the definition of $U'_C$ in the new representation, $((\widehat C\psi)' = U'_C\widetilde{\bar\psi'})$, we find
\begin{equation}
U'_C = UU_C\widetilde U.
\tag{26.20}
\end{equation}
Transformation~(26.20) agrees with the transformation of the $\gamma$ matrices only when $U$ is real. Thus expression~(26.5) is valid only in representations obtained from the spinor or standard representation by a real transformation.

The matrix~(26.5) is unitary, and transposition changes its sign:
\begin{equation}
U_CU^+_C = 1,\qquad \widetilde U_C = -U_C.
\tag{26.21}
\end{equation}
These properties are invariant under~(26.20) and therefore hold in every representation. The matrix~(26.5) is also Hermitian ($U_C = U^+_C$), but in general this property is not preserved by~(26.20).

\SourcePage{118}{123}
All these statements, including~(26.21), apply equally to $U_T$.

In second-quantized theory, the transformations $C$, $P$, and $T$ of the $\psi$-operators must be formulated as transformation rules for particle creation and annihilation operators. These rules may be established (as in \S\,13 for spin-0 particles) by requiring that the transformed $\psi$-operators admit the representation
\begin{equation}
\begin{aligned}
\widehat\psi^C(t,\mathbf{r}) &= U_C\widetilde{\widehat{\bar\psi}}(t,\mathbf{r}),\\
\widehat\psi^P(t,\mathbf{r}) &= i\gamma^0\widehat\psi(t,-\mathbf{r}),\\
\widehat\psi^T(t,\mathbf{r}) &= U_T\widetilde{\widehat{\bar\psi}}(-t,\mathbf{r}).
\end{aligned}
\tag{26.22}
\end{equation}

\ProblemHeading

Find the charge-conjugation operator in the Majorana representation (see problem~2, \S\,21).

\SolutionHeading The matrix $U'_C$ in the Majorana representation is obtained from the standard-representation matrix $U_C = -\alpha_y$ by transformation~(26.20), with $U = (\alpha_y + \beta)/\sqrt{2}$, and equals $U'_C = \alpha_y$ (here $\alpha_y$ and $\beta$ denote the standard-representation matrices). Using primes for quantities in the Majorana representation, we have $\widehat C\psi' = U'_C(\psi'^*\beta')$; since $\beta' = \alpha_y$,
\[
\widehat C\psi' = \alpha_y(\psi'^*\alpha_y) = \alpha_y\widetilde\alpha_y\psi'^* = \psi'^*,
\]
so charge conjugation is equivalent to complex conjugation.
